% ---------------------------------------------------------------------------
% Workshop paper, Section 3 (Method). Condensed from thesis Chapter 4.
% Cut relative to the thesis: in-context demonstrations (-> appendix), CORN,
% the KL anchor, task-description dropout, label noise, the LoRA loss study,
% and both in-line figures (-> appendix).
% ---------------------------------------------------------------------------
\section{Method}
\label{sec:method}

We keep Robometer's backbone and its per-frame design and change two things: the
data the model learns from, and the loss it minimizes.

\subsection{A failure-annotated dataset}
\label{sec:method-data}

Robometer trains its two deployment heads, progress and success, only on
successful trajectories, where an expert advancing steadily toward the goal hands
each frame a target for free: progress at frame $t$ is simply $t/T$. Failures
reach the model only indirectly, as the lower-ranked side of a preference
comparison, so neither deployment head ever sees failure as a direct target. The
success head is easy to fix, since a failed attempt is incomplete at every frame.
The progress head is not: a failure does not advance steadily, so its progress
cannot be read off the frame index. What is missing is a corpus whose failures
carry \emph{dense per-frame} progress. We build one by re-curating and expanding
RBM-1M, the corpus behind Robometer.

\paragraph{Labeling failures densely.}
Every failure frame receives a rubric score $s \in \{1,\dots,5\}$, mapped to a
normalized progress value $p = (s-1)/4$. A genuine failure never completes the
task, so its scores stay in $\{1,\dots,4\}$ and $p$ never reaches $1$. Three
routes produce these scores, chosen by data source. (i)~For failures we
\emph{generate} in simulation, MetaWorld \citep{yu_meta-world_2019,
mclean_meta-world_2025} and FailSafe \citep{lin_failsafe_2025}, progress comes
straight from simulator state, with no model in the loop. We pick a target score,
run a scripted expert until a task-specific injection point, then switch to a
noise mode that induces a particular failure (dropping a grasped object, stalling
a mechanism, pushing sideways) and keep the episode only if the state-based
detector confirms it ended at the intended score. (ii)~For failures where we have
only video, the real-world clips in RBM-1M and the DROID
\citep{khazatsky_droid_2024} failures we add, we infer the score with a two-stage
pipeline: a vision--language model describes each frame \emph{in isolation}, then
a language model, which never sees pixels, grades the descriptions cumulatively.
Splitting the work this way targets where these models actually fail. Asked to
label a whole clip at once, a VLM infers rather than observes: it reads the task
description, assumes the robot carried it out, and reports progress that never
happened. (iii)~RoboReward \citep{roboreward} partial executions are successes cut
short, hence monotone by construction, and are labeled by an even ramp up to the
clip's own score.

Because a per-frame rubric would give failures a stair-step target while
successes get a smooth ramp, we interpolate each failure linearly between the
frames where its level changes. Failures and successes then share one continuous
target on $[0,1]$, so the model cannot separate them by the \emph{shape} of the
target alone.

\paragraph{Composition.}
Table~\ref{tab:data} reports the mix. The corpus holds $648{,}714$ clips, of
which $116{,}786$ are failures. Crucially, the failures Robometer already held
become direct training signal for both deployment heads for the first time.
Appendix~\ref{apx:data} gives per-source detail and the labeling prompts.

\begin{table}[t]
\centering
\small
\caption{Composition of the training corpus. A row is one video clip rendered
from one camera view. \emph{Labels} names the source of the dense per-frame
failure targets; successes everywhere use the $t/T$ heuristic.}
\label{tab:data}
\begin{tabular}{@{}llrr@{}}
\toprule
Source & Labels & Failures & Successes \\
\midrule
Robometer expert demos      & ---                & $0$      & $490{,}798$ \\
Robometer mixed archives    & VLM+LLM            & $68{,}923$ & $26{,}029$ \\
MetaWorld (ours)            & simulator          & $24{,}402$ & $1{,}850$ \\
FailSafe (ours)             & simulator          & $2{,}098$  & $75$ \\
DROID (ours)                & VLM+LLM            & $5{,}503$  & $5{,}863$ \\
RoboReward (ours)           & ramp to score      & $15{,}860$ & $7{,}313$ \\
\midrule
Total                       &                    & $116{,}786$ & $531{,}928$ \\
\bottomrule
\end{tabular}
\end{table}

\subsection{Model}
\label{sec:method-arch}

We make two changes to Robometer's architecture. First, we drop the preference
head and keep only progress and success. The preference head existed to pull
failure signal into training indirectly; our dataset removes that constraint, so
the comparison route supplies nothing the two task heads do not already get. The
objective becomes $\mathcal{L} = \mathcal{L}_{\mathrm{prog}} +
\mathcal{L}_{\mathrm{succ}}$, and the two heads that remain are exactly the two
used at deployment. Second, we sample sixteen frames per clip instead of eight.
Reading a failure well means locating the moment it turns, the instant a grasp
slips, and a finer grain lets a per-frame label sit closer to that moment.

\subsection{An asymmetric objective}
\label{sec:method-loss}

Fine-tuned reward models over-predict success, reading completion into
trajectories that never finish. Robometer's loss treats the two ways of being
wrong as equals, but once a policy optimizes against the reward they are not
equally harmful \citep{gumbsch_demonstrations_2026}. A false positive plants a
spurious goal the policy will seek out and exploit; a false negative only slows
learning and offers no shortcut. The reward model should therefore be
conservative, and we make it so by weighting over-prediction more heavily than
under-prediction, so the model settles \emph{below} its targets.

On the progress head we compare the reported progress, the mean of the predicted
distribution $\bar{p}_t = \sum_i \hat{p}_{t,i} c_i$, against the target $p_t$. We
keep the C51 cross-entropy and multiply it by a per-frame weight,
\begin{equation}
\mathcal{L}_{\mathrm{prog}} = \frac{1}{T}\sum_{t=1}^{T} w_t\,
  \mathrm{CE}\!\big(\mathrm{Proj}(p_t),\, \hat{p}_t\big), \qquad
w_t = \begin{cases} 1, & \bar{p}_t > p_t,\\[2pt] \lambda, & \bar{p}_t \le p_t. \end{cases}
\label{eq:prog-asym}
\end{equation}
The comparison of $\bar{p}_t$ against $p_t$ reads the prediction as a fixed
reference, so it sets each frame's weight without contributing a gradient.

On the success head over-prediction means one specific thing, calling a frame
complete when it is not, so the bias falls on the label classes rather than on a
per-frame comparison. Incomplete frames keep full weight and complete frames are
damped:
\begin{equation}
\mathcal{L}_{\mathrm{succ}} =
  \frac{1}{|\mathcal{N}|}\sum_{t \in \mathcal{N}} \mathrm{BCE}\!\big(0, \hat{s}_t\big)
  \;+\; \lambda\,\frac{1}{|\mathcal{P}|}\sum_{t \in \mathcal{P}} \mathrm{BCE}\!\big(1, \hat{s}_t\big),
\label{eq:succ-asym}
\end{equation}
with $\mathcal{N} = \{t : s_t = 0\}$ and $\mathcal{P} = \{t : s_t = 1\}$. Damping
the positive class makes the model reluctant to fire a success signal. A failed
trajectory sits entirely in $\mathcal{N}$, so the damped term is supplied by the
successful trajectories it is batched with. We use $\lambda = 0.3$ on both heads;
$\lambda = 1$ recovers the symmetric loss. Appendix~\ref{apx:lambda} reports the
sweep behind this value.

\subsection{Training and model variants}
\label{sec:method-training}

All models are full fine-tunes of the language model and the reward heads, with
the vision encoder frozen, sharded across eight H100 GPUs with FSDP in bfloat16.
We vary two things.

\paragraph{Starting point.} We train from Robometer-4B, loading the released
reward model with its pretrained heads, and from the base vision--language model
Qwen3.5-4B \citep{team_qwen35-omni_2026} with freshly initialized heads and no
reward pretraining. The contrast isolates how much Robometer's existing reward
pretraining is worth. Because its heads start cold, the Qwen3.5 model runs in two
phases, a first pass over a broader corpus (including held-out humanoid and
human-only data) to settle the heads, then a second on the same standard-arm pool
the Robometer-4B runs use.

\paragraph{Ablation.} For each starting point we train two versions that switch
on our contributions one at a time. The first keeps Robometer's original
symmetric losses, isolating what the re-curated failure data alone contributes.
The second replaces them with the asymmetric objective of
Section~\ref{sec:method-loss}. Two starting points and two versions give four
models, evaluated against each other and against the untuned Robometer baseline.

\paragraph{Optimization.} AdamW \citep{loshchilov_decoupled_2019}, weight decay
$0.01$, cosine schedule with $10\%$ warmup, gradients clipped to norm $10$, eight
trajectories per GPU. Robometer-4B runs use a learning rate of $10^{-5}$ for
$5{,}000$ steps; Qwen3.5 runs use $2\times10^{-5}$ for $6{,}500$ steps in phase
one and $10^{-5}$ for $5{,}000$ in phase two. Rather than keep the corpus's
success-heavy mix, we oversample failures roughly twofold, raising their share of
what the model sees from $18\%$ to about $30\%$.
